\documentclass[../main.tex]{subfiles}

\graphicspath{{../image/}{../figure/}}

\begin{document}

\chapter{卷积网络}\label{chapter:ch10-卷积网络}

最简单的机器学习模型假设观测数据值是非结构化的，这意味着我们事先对数据向量$\mathbf{x}=(x_1,\cdots,x_D)$中各个元素之间可能存在的关系一无所知。如果我们对这些变量进行随机排列，并在所有训练数据和测试数据上一致地应用这种固定排列，那么到目前为止我们所考虑过的模型的性能将不会有任何差异。

然而，机器学习的许多应用涉及结构化数据（structured data），其中输入变量之间存在额外的关系。例如，自然语言中的单词形成一个序列（sequence）（参见\cref{chapter:ch12-Transformer与注意力机制}）：如果我们将语言建模为一个生成式自回归过程，那么我们将期望每个单词与直接邻近的单词更相关，而不是依赖于序列中更早出现的单词。同样，图像的像素之间具有良定的空间关系，其中输入变量被安排在二维网格中，相邻像素具有高度相关的值。

我们已经可以通过在训练目标的误差函数中添加正则化项（参见\cref{chapter:ch09-正则化}），以及通过数据增强或修改模型架构来利用我们对特定数据模态结构的知识。这些方法可以帮助引导模型尊重输入数据在变换方面的某些属性，如不变性和等变性（参见\cref{subsection:ch09-对称性不变性}）。在本章中，我们将介绍一种称为卷积神经网络（Convolutional Neural Network, CNN）的架构方法。CNN 可以视为具有参数共享的稀疏连接多层网络，并用于编码特定于图像数据的不变性和等变性。

\subfile{Chapter-10/section-01.tex}
\subfile{Chapter-10/section-02.tex}
\subfile{Chapter-10/section-03.tex}
\subfile{Chapter-10/section-04.tex}
\subfile{Chapter-10/section-05.tex}
\subfile{Chapter-10/section-06.tex}
\subfile{Chapter-10/section-07.tex}

\end{document}
